\documentclass[11pt]{article}

\newcommand{\cnum}{Math 132H}
\newcommand{\ced}{Winter 2026}
\newcommand{\ctitle}[4]{\title{\vspace{-0.5in}\cnum, \ced\\Assignment #1: #2}\author{\vspace{-0.35in}\\#3, #4}}
\usepackage{enumitem}
\usepackage[usenames,dvipsnames,svgnames,table,hyperref]{xcolor}
\usepackage{amsmath, amsfonts}
\usepackage{graphicx} % Required for inserting images
\usepackage{float}
\usepackage{listings}
\usepackage{xcolor}
\usepackage{hyperref}
\usepackage{comment}

\renewcommand*{\theenumi}{\alph{enumi}}
\renewcommand*\labelenumi{(\theenumi)}
\renewcommand*{\theenumii}{\roman{enumii}}
\renewcommand*\labelenumii{\theenumii.}

\definecolor{codegreen}{rgb}{0,0.6,0}
\definecolor{codegray}{rgb}{0.5,0.5,0.5}
\definecolor{codepurple}{rgb}{0.58,0,0.82}
\definecolor{backcolour}{rgb}{0.95,0.95,0.92}
\definecolor{header}{HTML}{205459}
\definecolor{solution}{HTML}{204d52}

\newcommand{\solution}[1]{{{\textcolor{header}{Solution:} \textcolor{solution}{#1}}}}
\setlist[enumerate]{label=(\alph*)}

\lstdefinestyle{mystyle}{
    backgroundcolor=\color{backcolour},
    commentstyle=\color{codegreen},
    keywordstyle=\color{magenta},
    numberstyle=\tiny\color{codegray},
    stringstyle=\color{codepurple},
    basicstyle=\ttfamily\footnotesize,
    breakatwhitespace=false,
    breaklines=true,
    captionpos=b,
    keepspaces=true,
    numbers=left,
    numbersep=5pt,
    showspaces=false,
    showstringspaces=false,
    showtabs=false,
    tabsize=2
}


\begin{document}
\ctitle{\#1}{}{Asher Christian}{006-150-286}
\date{}
\maketitle
\section{Exercise 1.1}
\emph{
    Describe geometrically the sets of points $z$ in the complex plane
}
\solution{
    \begin{enumerate}
        \item $|z-z_1| = |z-z_2|$ A straight line passing through $(z_1+z_2)$
        \item $\frac{1}{z} = \overline{z}$. If $z = a + ib$ then  $\frac{1}{z} = \frac{a-ib}{a^2 + b^2}$ and $\overline{ z} = a-ib$.
            Setting the two equal we see that the $z$ that satisfy this are all those such that $a^2 + b^2 = 1$ or the ones htat sit on the unit circle of radius one. The points
            $z = e^{i\theta}$ for some $\theta$ real number.
        \item $\Re(z) = 3$ the line going straight up that passes through  $z = 3$ and includes all  $z$ of form $z = 3 +ib$ for any  $b$ real.
        \item  $\Re(z) > c$ the open half plane to the right whose border is defined as the line  $\Re(z) = c$ as before.
        \item $\Re(az + b) > 0$ where  $a,b \in \mathbb{C}$ if we look at $a = re^{i\theta}$ then this is taking the half plane $\Re(z) > 0$, rotating it by angle $\theta$ and then shifting it
            by the real part of $b$.
        \item  $|z| = \Re(z) + 1$ rewriting  $z = a +ib$ we get that this is the equation
             $a = \frac{b^2}{2} -\frac{1}{2}$ which we recognize as a quadratic equation, so this is a parabola.
         \item $\Im(z) = c$ with  $c \in \mathbb{R}$. A horizontal line in the plane that passes through $z = ic$
    \end{enumerate}
}
\section{Exercise 1.2}
\emph{
    let $\langle \cdot , \cdot \rangle$ denote the usual inner product in $ \mathbb{R}^2$.
    Define a Hermitian inner product $(\cdot, \cdot)$ in  $ \mathbb{C}$ by 
    \[
        (z,w) = z \overline{w}
    \] 
    Show that
    \[
        \langle z, w \rangle = \frac{1}{2}\left[(z,w) + (w,z)\right] = \Re(z,w)
    \] 
}
\solution{
    If $z = x_1 + iy_1$ and $w = x_2 + iy_2$ then
    \[
    \langle z, w \rangle = x_1x_2 + y_1y_2
    \] 
    and
    \[
        (z,w) = x_1x_2  + y_1y_2 -i(x_1y_2 -y_1x_2) \qquad (w,z) = x_1x_2 + y_1y_2 -i(y_1x_2 - y_2x_1)
    \] 
    adding them together we get
    \[
        \frac{1}{2}[(z,w) + (w,z)] = x_1x_2 + y_1y_2 = \Re(z,w) = \langle z, w \rangle
    \] 
}

\section{Exercise 1.4}
\emph{
    Show that it is impossible to define a total ordering on $ \mathbb{C}$. One cannot find a relation $\succ$ between complex numbers
    so that
     \begin{enumerate}
         \item For any two complex numbers $z,w$ one and only one of the following is true:  $z \succ w, w\succ z$ or  $z = w$. 
         \item For all $z_1,z_2,z_3 \in \mathbb{C}$ the relation $z_1 \succ z_2$ implies $z_1 + z_3 \succ z_2 + z_3$
         \item for all $z_1,z_2,z_3 \in \mathbb{C}$ with $z_3 \succ 0$ then $z_1 \succ z_2$ implies $z_1z_3 \succ z_2z_3$
    \end{enumerate}
}
\solution{
    Assume for contradiction that such a $\succ$ relation exists. Assume further that  $i \succ 0$. Then
    we get by $(iii)$
     \[
         i^2 \succ i 0 \implies -1 \succ 0
    \] 
    repeating further we get
    \[
    -i \succ 0 \qquad 1 \succ 0
    \] 
    however
    \begin{align*}
        -1 \succ 0\\
        -1 + 1 \succ 1\\
        0 \succ 1
    \end{align*}
    this contradicts $1 \succ 0$ so our second assumption must be false. Therefore if a relation exists
    it must satisfy  $0 \succ i$. It also must satisfy  $0 \succ -1$ for if we assume  $-1 \succ 0$ then we can say $-1 + 1 \succ 0 + 1$ which implies $0 \succ 1$ and
    we can also say that  $-1 \cdot -1 \succ -1 * 0$ so  $1 \succ 0$ which is a contradiction. Thus we must have both  $0 \succ i$ and $0 \succ -1$ and  $1 \succ 0$ and $-i \succ 0$.
    However
    \[
    -i \succ 0 \implies 0 * -i \succ i * -i \implies 0 \succ 1
    \] 
    a contradiction. Thus it is impossible to define a total ordering on $ \mathbb{C}$
}

\section{Exercise 1.7}
\emph{
    \begin{enumerate}
        \item Let $z,w$ be two complex numbers such that  $\overline{z}w \ne 1$. Prove that
            \[
                \left | \frac{w-z}{1-\overline{w}z} \right | < 1 \qquad \text{if $|z| < 1$ and  $|w| < 1$}
            \] 
            and also that
            \[
                \left | \frac{w-z}{1- \overline{w}z} \right| = 1 \qquad \text{if $|z| = 1$ or  $|w| = 1$}
            \] 
        \item Prove that for a fixed $w$ in the unit disc $ \mathbb{D}$, the mapping
            \[
            F : z \rightarrow \frac{w-z}{1-\overline{w}z}
            \] 
            satisfies
            \begin{enumerate}
                \item $F$ maps the unit disc to itself and is holomorphic
                \item $F$ interchanges $0$ and $w$ 
                \item $|F(z)| = 1$ if  $|z| = 1$
                \item  $F : \mathbb{D} \rightarrow \mathbb{D}$ is bijective
            \end{enumerate}
    \end{enumerate}
}
\solution{
    let $w = r_2e^{i\theta_2}$ and $z = r_1e^{i\theta_1}$ then we can write it as
    \[
    \left | \frac{r_2e^{i\theta_2} - r_1e^{i\theta_1}}{1 - r_1r_2e^{i(\theta_1-\theta_2)}} \right | = \left| e^{i\theta_1} \frac{r_2e^{i(\theta_2-\theta_1)} - r_1}{1 - r_1r_2e^{-i(\theta_2-\theta_1)}} \right|
    \] 
    and since rotation does not affect magnitude we can get rid of the $e^{i\theta_1}$ without affecting the total. Additionalyl $\tilde w = r_2e^{i(\theta_2 - \theta_1)}$ is in the unit disc since its magnitude is unchanged. So we can assume that $z$ is real
    so it now suffices to show for all  $|r| \le 1$ 
     \[
    |w-r| \le |1-\overline{w}r|
    \] 
    or, squaring both sides equivalently
    \[
        (r-w)(r-\overline{w}) \le (1-rw)(1-r\overline{w})
    \] 
    but
    \[
        (r-w)(r-\overline{w}) = r^2 - r \overline{w} - rw + w \overline{w} = r^2 - r(w + \overline{w}) + |w|^2
    \] 
    and
    \[
        (1-rw)(1-\overline{rw}) = 1 - r \overline{w} - rw + r^2 w \overline{w} = 1 - r(w + \overline{w}) + r^2 |w|^2
    \] 
    we want to show
    \[
    r^2 - r(w + \overline{w}) + |w|^2 \le 1 - r(w + \overline{w}) + r^2 |w|^2
    \] 
    or equivalently
    \[
    r^2 + |w|^2 \le 1 + r^2|w|^2 
    \] 
    and moving it all to one side
    \[
    0 \le 1 + r^2|w|^2 - r^2 - |w|^2 = (1-r^2) - |w|^2(1-r)^2 = (1-r^2)(1-|w|^2)
    \] 
    but since we assumed that $r \le 1$ and  $|w| \le 1$ we see that the inequality holds with equality holding only if  $|r| = 1$ or  $|w| = 1$ as given.\\
    For part $(b)$ Fix  $w$ in the unit disc. then since  $|F(z)| \le 1$ we have that the image of  $F$ under  $z \in \mathbb{D}$ is in  $\mathbb{D}$.
    Additionally $F$ is holomorphic since it is the ratio of two holomorphic functions. It is holomorphic on the set $1 - \overline{w}z \ne 0$. which is satisfied by $z = w$ if  $|w| =1$ and is not satisfied within the unit disc otherwise.
    Additionally $F(0) = w$ and  $F(w) = 0$ trivially. As proven in the previous part  $|F(z)| = 1$ if  $|z| = 1$. Assume now that  $|w| < 1$ then consider
     \[
    F \circ F = id
    \] 
    manually check
}

\section{Exercise 1.8}
\emph{
    Suppose $U$ and $V$ are open sets in the complex plane. Prove that if $f : U \rightarrow V$ and $g : V \rightarrow \mathbb{C}$
    are two functions that are differentiable as functions of two real variables, and $h = g \circ f$ then
     \[
    \frac{\partial h}{\partial z} = \frac{\partial g}{\partial z} \frac{\partial f}{\partial z} + \frac{\partial g}{\partial \overline{z}} \frac{\partial \overline{f}}{\partial z}
    \] 
    and
    \[
        \frac{\partial h}{\partial \overline{z}} = \frac{\partial g}{\partial z} \frac{\partial f}{\partial \overline{ z}} + \frac{\partial g}{\partial \overline{ z}} \frac{\partial \overline{ f}}{\partial \overline{ z}}
    \] 
}
\solution{
    We have
    \[
        \frac{\partial h}{\partial z} = \frac{1}{2}(\frac{\partial h}{\partial x} + \frac{1}{i}\frac{\partial h}{\partial y})
    \] 
    where
    \[
        \frac{\partial h}{\partial x}(x,y) = \lim_{r\to 0} \frac{h(x+r,y)-h(x,y)}{r}
    \] 
}
\end{document}

\section{Exercise 1.9}
\emph{
    Show that in polar coordinates, the Cauchy-Riemann equations take the form
    \[
        \frac{\partial u}{\partial r} = \frac{1}{r} \frac{\partial v}{\partial \theta} \qquad \frac{1}{r} \frac{\partial u}{\partial \theta} = -\frac{\partial v}{\partial r}
    \] 
    Use these equations to show that the logarithm function defined by
    \[
        \log z = \log r + i\theta \qquad z = re^{i\theta} \text{ with } -\pi < \theta < \pi
    \] 
    is holomorphic in the region $r > 0$ and  $-\pi < \theta < \pi$
}
